% Gerado por docs/sistema/gera_modelo_tex.py a partir de docs/sistema/modelo.json. Não editar à mão.

\chapter{Dicionário de dados}
\label{ap:dicionario}

Gerado a partir do modelo do back-end. Cada tabela traz as colunas, o tipo, se aceita nulo, as chaves e a descrição. As relações estão na Tabela~\ref{tab:relacoes}.

\section*{auditoria}
Trilha de auditoria somente de acréscimo, encadeada por HMAC-SHA256, protegida por gatilhos contra UPDATE, DELETE e TRUNCATE; registra cada ação no sistema (RNF02).

\begin{small}
\begin{longtable}{p{3.4cm}p{1.9cm}p{0.8cm}p{1.6cm}p{6.6cm}}
\hline
\textbf{Coluna} & \textbf{Tipo} & \textbf{Nulo} & \textbf{Chave} & \textbf{Descrição} \\ \hline\endhead
seq & inteiro & não & PK & Número sequencial da entrada, gerado pelo banco; ordena a cadeia e entra no cálculo do hash. \\
versao & inteiro & não &  & Versão do formato da cadeia. O código só aceita 2 (HMAC-SHA256 com chave fora do banco). \\
quando & texto(40) & não &  & Data e hora da ação, ISO 8601 em UTC com sufixo Z. \\
usuario & texto(100) & não &  & Login de quem executou a ação; o agendador interno grava 'agendador'. \\
acao & texto(60) & não &  & Nome da ação no formato objeto.ação, por exemplo caso.ver, coleta.executar, alerta.decidir, caso.acesso\_negado. \\
caso\_id & texto(32) & sim &  & ID do caso afetado (aponta para casos, sem chave estrangeira); nulo em ações fora de caso, como login e logout. \\
justificativa & texto & sim &  & Justificativa informada na requisição; nula em ações sem justificativa, como login e logout. \\
detalhes & JSON & não &  & JSON com os dados da ação (inclui o perfil do usuário). Textos longos e objetos grandes são substituídos por tamanho e SHA-256. \\
hash\_anterior & texto(64) & não &  & Hash da entrada anterior, em hexadecimal; a primeira entrada usa 64 zeros. \\
hash & texto(64) & não &  & HMAC-SHA256 do JSON canônico da entrada sem o hash, com chave fora do banco; o topo é copiado para uma âncora externa. \\
\hline
\end{longtable}
\end{small}

\section*{casos}
Caso de investigação: agrupa coletas, entidades, alertas e equipe, com finalidade declarada e prazo de retenção (RNF04).

\begin{small}
\begin{longtable}{p{3.4cm}p{1.9cm}p{0.8cm}p{1.6cm}p{6.6cm}}
\hline
\textbf{Coluna} & \textbf{Tipo} & \textbf{Nulo} & \textbf{Chave} & \textbf{Descrição} \\ \hline\endhead
id & texto(32) & não & PK & Identificador do caso, UUID4 em hexadecimal (32 caracteres), gerado pelo sistema. \\
titulo & texto(300) & não &  & Título do caso, informado na criação. \\
descricao & texto & não &  & Descrição livre do caso, informada na criação; pode ficar vazia. \\
numero\_procedimento & texto(100) & não &  & Número do procedimento administrativo ou judicial, informado na criação; até 100 caracteres, pode ficar vazio. \\
criado\_por & texto(100) & não &  & Login de quem criou o caso; na criação essa pessoa entra em caso\_membros como responsável. \\
criado\_em & texto(40) & não &  & Data e hora de criação, ISO 8601 em UTC. \\
finalidade & texto & não &  & Finalidade declarada do tratamento dos dados, obrigatória na criação; vazia em casos anteriores à coluna (RNF04). \\
retencao\_ate & texto(10) & sim &  & Prazo de retenção, data AAAA-MM-DD inclusive. Depois dele o caso fica só leitura; nulo em casos anteriores à coluna. \\
\hline
\end{longtable}
\end{small}

\section*{usuarios}
Conta de usuário real do sistema, com senha em hash Argon2id, segundo fator TOTP cifrado e controle de bloqueio (RNF01).

\begin{small}
\begin{longtable}{p{3.4cm}p{1.9cm}p{0.8cm}p{1.6cm}p{6.6cm}}
\hline
\textbf{Coluna} & \textbf{Tipo} & \textbf{Nulo} & \textbf{Chave} & \textbf{Descrição} \\ \hline\endhead
id & texto(32) & não & PK & Identificador da conta, UUID4 em hexadecimal, gerado pelo sistema. \\
usuario & texto(100) & não & única & Login, único, em minúsculas; 3 a 64 caracteres entre letras, dígitos, ponto, hífen e sublinhado. \\
nome & texto(200) & não &  & Nome da pessoa titular da conta; pode ficar vazio. \\
perfil & texto(20) & não &  & Perfil de acesso global. Valores: analista, delegado, admin. \\
senha\_hash & texto(300) & não &  & Hash Argon2id da senha (biblioteca argon2-cffi); a senha em claro não é guardada. \\
ativo & lógico & não &  & Indica se a conta pode entrar no sistema; conta inativa tem as sessões encerradas. \\
mfa\_segredo\_cifrado & texto & sim &  & Segredo TOTP cifrado com AES-256-GCM, em Base64, chave fora do banco; nulo até o cadastro ou após redefinição. \\
mfa\_ativo & lógico & não &  & Indica se o segundo fator já foi confirmado; passa a verdadeiro no primeiro código TOTP aceito. \\
mfa\_confirmado\_em & texto(40) & sim &  & Data e hora em que o segundo fator foi confirmado, ISO 8601 em UTC; nulo antes disso. \\
mfa\_ultimo\_passo & inteiro & sim &  & Último passo de tempo TOTP aceito; impede reutilizar o mesmo código (RFC 6238, seção 5.2). \\
falhas\_seguidas & inteiro & não &  & Contador de falhas de senha ou de código seguidas; zera no acerto ou ao atingir o limite e bloquear. \\
bloqueado\_ate & texto(40) & sim &  & Data e hora até a qual a conta fica bloqueada após falhas seguidas, ISO 8601 em UTC; nulo sem bloqueio. \\
criado\_por & texto(100) & não &  & Login de quem criou a conta; pode ficar vazio. \\
criado\_em & texto(40) & não &  & Data e hora de criação da conta, ISO 8601 em UTC. \\
senha\_alterada\_em & texto(40) & não &  & Data e hora da última definição da senha, ISO 8601 em UTC. \\
\hline
\end{longtable}
\end{small}

\section*{alertas}
Alerta gerado pelas tipologias sobre entidades do caso, com evidências, severidade e a última decisão do investigador (RF12, RF14).

\begin{small}
\begin{longtable}{p{3.4cm}p{1.9cm}p{0.8cm}p{1.6cm}p{6.6cm}}
\hline
\textbf{Coluna} & \textbf{Tipo} & \textbf{Nulo} & \textbf{Chave} & \textbf{Descrição} \\ \hline\endhead
id & texto(32) & não & PK & Identificador do alerta, UUID4 em hexadecimal; mantido quando a reexecução detecta o mesmo alerta. \\
caso\_id & texto(32) & não & FK casos & Caso a que o alerta pertence (chave estrangeira para casos). \\
tipologia & texto(5) & não &  & Código da tipologia que gerou o alerta. Valores: T1 a T10. \\
titulo & texto(300) & não &  & Título do alerta, gerado pela regra da tipologia. \\
motivo & texto & não &  & Explicação, gerada pela regra, de por que o alerta disparou. \\
severidade & texto(10) & não &  & Severidade definida pela tipologia na tabela de políticas, não pelo alerta. Valores: alta, media, baixa. \\
entidades & JSON & não &  & Lista JSON de IDs de entidades envolvidas (aponta para entidades, sem chave estrangeira). \\
evidencias & JSON & não &  & Lista JSON de evidências: descrição, fonte, captura\_id, url\_ou\_arquivo e sha256 da captura citada (aponta para capturas, sem chave estrangeira). \\
parametros & JSON & não &  & JSON com os valores observados pela regra e, em 'limiares', os limiares numéricos usados na execução. \\
criado\_em & texto(40) & não &  & Data e hora em que o alerta foi criado, ISO 8601 em UTC. \\
chave & texto(64) & não &  & SHA-256 de tipologia, IDs das entidades e discriminador; identidade estável do alerta entre execuções das tipologias. \\
ativo & lógico & não &  & Falso quando a última execução não detectou mais o alerta, mas ele foi mantido por ter decisão registrada. \\
situacao & texto(12) & não &  & Situação atual, cópia da última decisão. Valores: pendente, confirmado, descartado. \\
decidido\_por & texto(100) & sim &  & Login de quem tomou a última decisão; nulo enquanto pendente. \\
decidido\_em & texto(40) & sim &  & Data e hora da última decisão, ISO 8601 em UTC; nulo enquanto pendente. \\
justificativa\_decisao & texto & sim &  & Justificativa da última decisão; o histórico completo fica em alerta\_decisoes. \\
\hline
\end{longtable}
\end{small}

\section*{anotacoes}
Anotação de um membro sobre uma entidade ou um alerta do caso; só de acréscimo, protegida por gatilhos contra alteração e exclusão (RF08).

\begin{small}
\begin{longtable}{p{3.4cm}p{1.9cm}p{0.8cm}p{1.6cm}p{6.6cm}}
\hline
\textbf{Coluna} & \textbf{Tipo} & \textbf{Nulo} & \textbf{Chave} & \textbf{Descrição} \\ \hline\endhead
id & texto(32) & não & PK & Identificador da anotação, UUID4 em hexadecimal. \\
caso\_id & texto(32) & não & FK casos & Caso a que a anotação pertence (chave estrangeira para casos). \\
alvo\_tipo & texto(10) & não &  & Tipo do objeto anotado. Valores: entidade, alerta. \\
alvo\_id & texto(32) & não &  & ID do objeto anotado; aponta para entidades ou alertas conforme alvo\_tipo, sem chave estrangeira. \\
autor & texto(100) & não &  & Login de quem escreveu a anotação. \\
criado\_em & texto(40) & não &  & Data e hora da anotação, ISO 8601 em UTC. \\
texto & texto & não &  & Texto da anotação. Não se edita: corrigir é criar outra anotação. \\
\hline
\end{longtable}
\end{small}

\section*{caso\_membros}
Membros de cada caso e seu papel; controla quem pode ver e agir no caso (RNF01).

\begin{small}
\begin{longtable}{p{3.4cm}p{1.9cm}p{0.8cm}p{1.6cm}p{6.6cm}}
\hline
\textbf{Coluna} & \textbf{Tipo} & \textbf{Nulo} & \textbf{Chave} & \textbf{Descrição} \\ \hline\endhead
id & inteiro & não & PK & Identificador numérico sequencial, gerado pelo banco. \\
caso\_id & texto(32) & não & FK casos & Caso a que o membro pertence (chave estrangeira para casos). \\
usuario & texto(100) & não &  & Login do membro, de conta real ou de token de desenvolvimento (TOKENS\_DEV); sem chave estrangeira para usuarios. \\
papel & texto(15) & não &  & Papel no caso, em ordem crescente de permissão. Valores: leitor, investigador, responsavel. \\
adicionado\_por & texto(100) & não &  & Login de quem incluiu o membro; 'migracao' para casos anteriores à tabela. \\
adicionado\_em & texto(40) & não &  & Data e hora da inclusão, ISO 8601 em UTC. \\
\hline
\end{longtable}
\end{small}

\section*{coletas}
Registro de cada execução de coleta de uma fonte num caso, bem-sucedida ou não; recoletas apontam para a coleta repetida (RNF12, RF07).

\begin{small}
\begin{longtable}{p{3.4cm}p{1.9cm}p{0.8cm}p{1.6cm}p{6.6cm}}
\hline
\textbf{Coluna} & \textbf{Tipo} & \textbf{Nulo} & \textbf{Chave} & \textbf{Descrição} \\ \hline\endhead
id & texto(32) & não & PK & Identificador da coleta, UUID4 em hexadecimal. \\
caso\_id & texto(32) & não & FK casos & Caso em que a coleta foi feita (chave estrangeira para casos). \\
fonte & texto(50) & não &  & Fonte consultada. Valores aceitos pela API: cnpj\_arquivo, pncp, portal\_transparencia, pncp\_termos, tse\_arquivo, edital\_arquivo. \\
parametros & JSON & não &  & JSON com os parâmetros da consulta enviados pelo usuário, validados por fonte antes da execução. \\
conector\_nome & texto(50) & não &  & Nome do conector que executou a coleta; vazio se a coleta falhou antes de identificá-lo. \\
conector\_versao & texto(20) & não &  & Versão do conector que executou a coleta; vazia se não identificada. \\
situacao & texto(15) & não &  & Situação da coleta. Valores: em\_andamento, sucesso, falha, nao\_configurado. \\
tentativas & inteiro & não &  & Maior número de tentativas que uma operação da coleta precisou; 0 quando nada foi tentado. \\
motivo & texto & não &  & Mensagem de falha ou de fonte não configurada; vazio em caso de sucesso. \\
registros & inteiro & não &  & Quantidade de registros devolvidos pelo conector. \\
avisos & JSON & não &  & Lista JSON de avisos textuais do conector (até 50, de até 300 caracteres); em nao\_configurado, só o motivo. \\
usuario & texto(100) & não &  & Login de quem pediu a coleta; 'agendador' nas recoletas automáticas do monitoramento. \\
iniciado\_em & texto(40) & não &  & Data e hora de início, ISO 8601 em UTC; a linha é gravada antes de qualquer consulta externa. \\
concluido\_em & texto(40) & sim &  & Data e hora de fechamento, ISO 8601 em UTC; nulo enquanto em andamento. \\
recoleta\_de & texto(32) & sim & FK coletas & Coleta que esta repete, para comparação (chave estrangeira para coletas); nulo na coleta original. \\
\hline
\end{longtable}
\end{small}

\section*{sessoes}
Sessão de login de uma conta real, em etapas até o segundo fator; o token só fica como SHA-256 (RNF01).

\begin{small}
\begin{longtable}{p{3.4cm}p{1.9cm}p{0.8cm}p{1.6cm}p{6.6cm}}
\hline
\textbf{Coluna} & \textbf{Tipo} & \textbf{Nulo} & \textbf{Chave} & \textbf{Descrição} \\ \hline\endhead
id & texto(32) & não & PK & Identificador da sessão, UUID4 em hexadecimal. \\
token\_resumo & texto(64) & não & única & SHA-256 do token opaco entregue ao cliente; único. O token em claro não é guardado. \\
conta\_id & texto(32) & não & FK usuarios & Conta dona da sessão (chave estrangeira para usuarios). \\
etapa & texto(15) & não &  & Etapa da autenticação. Valores: mfa (falta o código), cadastro\_mfa (falta cadastrar o fator), ativa (única que dá acesso às rotas). \\
criada\_em & texto(40) & não &  & Data e hora de abertura da sessão, ISO 8601 em UTC. \\
expira\_em & texto(40) & não &  & Data e hora de expiração absoluta, ISO 8601 em UTC; definida pela duração configurada para a etapa. \\
ultima\_atividade\_em & texto(40) & não &  & Última atividade registrada, ISO 8601 em UTC; usada para encerrar por inatividade. \\
encerrada\_em & texto(40) & sim &  & Data e hora de encerramento, ISO 8601 em UTC; nulo enquanto a sessão está aberta. \\
motivo\_encerramento & texto(40) & não &  & Motivo do encerramento; vazio com a sessão aberta. Valores: logout, expirada, inatividade, conta\_inativa, bloqueio, segundo\_fator\_confirmado, troca\_de\_senha, alteracao\_da\_conta. \\
\hline
\end{longtable}
\end{small}

\section*{tarefas}
Tarefa atribuída dentro de um caso, com o estado atual; cada mudança fica em tarefa\_eventos (RF08).

\begin{small}
\begin{longtable}{p{3.4cm}p{1.9cm}p{0.8cm}p{1.6cm}p{6.6cm}}
\hline
\textbf{Coluna} & \textbf{Tipo} & \textbf{Nulo} & \textbf{Chave} & \textbf{Descrição} \\ \hline\endhead
id & texto(32) & não & PK & Identificador da tarefa, UUID4 em hexadecimal. \\
caso\_id & texto(32) & não & FK casos & Caso a que a tarefa pertence (chave estrangeira para casos). \\
titulo & texto(300) & não &  & Título da tarefa, informado pelo usuário. \\
descricao & texto & não &  & Descrição da tarefa; pode ficar vazia. \\
responsavel & texto(100) & sim &  & Login do membro do caso encarregado da tarefa, em minúsculas; nulo quando sem responsável. \\
situacao & texto(15) & não &  & Situação da tarefa. Valores: aberta, em\_andamento, concluida, cancelada. \\
prazo & texto(10) & sim &  & Prazo da tarefa, data AAAA-MM-DD; nulo quando sem prazo. \\
criado\_por & texto(100) & não &  & Login de quem criou a tarefa. \\
criado\_em & texto(40) & não &  & Data e hora de criação, ISO 8601 em UTC. \\
atualizado\_em & texto(40) & não &  & Data e hora da última alteração, ISO 8601 em UTC. \\
\hline
\end{longtable}
\end{small}

\section*{alerta\_decisoes}
Histórico das decisões sobre alertas, uma linha por decisão; só de acréscimo, protegido por gatilhos contra alteração e exclusão (RF12).

\begin{small}
\begin{longtable}{p{3.4cm}p{1.9cm}p{0.8cm}p{1.6cm}p{6.6cm}}
\hline
\textbf{Coluna} & \textbf{Tipo} & \textbf{Nulo} & \textbf{Chave} & \textbf{Descrição} \\ \hline\endhead
id & inteiro & não & PK & Identificador numérico sequencial, gerado pelo banco. \\
alerta\_id & texto(32) & não & FK alertas & Alerta decidido (chave estrangeira para alertas). \\
caso\_id & texto(32) & não & FK casos & Caso do alerta (chave estrangeira para casos). \\
decisao & texto(12) & não &  & Decisão tomada. Valores: confirmado, descartado. \\
justificativa & texto & não &  & Justificativa da decisão, obrigatória, informada pelo investigador. \\
usuario & texto(100) & não &  & Login de quem decidiu. \\
perfil & texto(20) & não &  & Perfil de acesso de quem decidiu no momento da decisão (analista, delegado ou admin). \\
decidido\_em & texto(40) & não &  & Data e hora da decisão, ISO 8601 em UTC. \\
\hline
\end{longtable}
\end{small}

\section*{capturas}
Registro de custódia de cada carga bruta coletada, com SHA-256, tamanho e local do arquivo; o papel da aplicação só insere e lê (RF07).

\begin{small}
\begin{longtable}{p{3.4cm}p{1.9cm}p{0.8cm}p{1.6cm}p{6.6cm}}
\hline
\textbf{Coluna} & \textbf{Tipo} & \textbf{Nulo} & \textbf{Chave} & \textbf{Descrição} \\ \hline\endhead
id & texto(32) & não & PK & Identificador da captura, UUID4 em hexadecimal. \\
caso\_id & texto(32) & não & FK casos & Caso em que a carga foi coletada (chave estrangeira para casos). \\
fonte & texto(50) & não &  & Nome do conector que gravou a carga, por exemplo pncp, cnpj\_arquivo, tse\_arquivo. \\
url\_ou\_arquivo & texto & não &  & URL consultada ou nome do arquivo de origem; começa com simulado:// quando a carga é sintética. \\
coletado\_em & texto(40) & não &  & Data e hora da coleta da carga, ISO 8601 em UTC. \\
sha256 & texto(64) & não &  & SHA-256 do conteúdo original, calculado antes de cifrar; usado na verificação de integridade. \\
tamanho\_bytes & inteiro & não &  & Tamanho do conteúdo original, em bytes. \\
caminho\_relativo & texto(200) & não &  & Caminho na pasta de custódia: dois primeiros caracteres do SHA-256, barra, SHA-256; extensão .gcm quando cifrado com AES-256-GCM. \\
metadados & JSON & não &  & JSON com metadados do conector; sempre traz 'simulado' e, se cifrado, 'cifra\_em\_repouso' com algoritmo e chave\_id. \\
conector\_nome & texto(50) & não &  & Nome do conector que produziu a captura; vazio em capturas anteriores à coluna. \\
conector\_versao & texto(20) & não &  & Versão do conector que produziu a captura; vazia em capturas anteriores à coluna. \\
coleta\_id & texto(32) & sim & FK coletas & Coleta a que a captura pertence (chave estrangeira para coletas); nulo em capturas anteriores à coluna. \\
\hline
\end{longtable}
\end{small}

\section*{tarefa\_eventos}
Histórico das tarefas, uma linha por criação ou alteração; só de acréscimo, protegido por gatilhos contra alteração e exclusão (RF08).

\begin{small}
\begin{longtable}{p{3.4cm}p{1.9cm}p{0.8cm}p{1.6cm}p{6.6cm}}
\hline
\textbf{Coluna} & \textbf{Tipo} & \textbf{Nulo} & \textbf{Chave} & \textbf{Descrição} \\ \hline\endhead
id & inteiro & não & PK & Identificador numérico sequencial, gerado pelo banco. \\
tarefa\_id & texto(32) & não & FK tarefas & Tarefa alterada (chave estrangeira para tarefas). \\
caso\_id & texto(32) & não & FK casos & Caso da tarefa (chave estrangeira para casos). \\
usuario & texto(100) & não &  & Login de quem criou ou alterou a tarefa. \\
quando & texto(40) & não &  & Data e hora do evento, ISO 8601 em UTC. \\
alteracoes & JSON & não &  & JSON {campo: {de, para}}. Na criação, 'de' é nulo e registra título, responsável, situação e prazo; mudança de descrição guarda só a contagem de caracteres. \\
\hline
\end{longtable}
\end{small}

\section*{entidades}
Entidade resolvida a partir das capturas do caso (empresa, pessoa, órgão etc.), única por caso, tipo e chave de resolução (RF02).

\begin{small}
\begin{longtable}{p{3.4cm}p{1.9cm}p{0.8cm}p{1.6cm}p{6.6cm}}
\hline
\textbf{Coluna} & \textbf{Tipo} & \textbf{Nulo} & \textbf{Chave} & \textbf{Descrição} \\ \hline\endhead
id & texto(32) & não & PK & Identificador da entidade, UUID4 em hexadecimal. \\
caso\_id & texto(32) & não & FK casos & Caso a que a entidade pertence (chave estrangeira para casos). \\
tipo & texto(20) & não &  & Tipo da entidade. Valores: empresa, pessoa, orgao, contrato, sancao, candidato, edital. \\
chave & texto(300) & não &  & Chave de resolução por tipo, por exemplo CNPJ básico para empresa, numeroControlePNCP para contrato; pessoa usa CPF descaracterizado e nome normalizado. \\
rotulo & texto(300) & não &  & Nome de exibição da entidade, como razão social ou nome da pessoa. \\
identificador & texto(100) & não &  & Identificador oficial exibido, como CNPJ, CPF descaracterizado ou número de controle; pode ficar vazio. \\
atributos & JSON & não &  & JSON com os atributos vindos das fontes; 'origem\_atributos' indica a captura que forneceu cada valor. CPF só descaracterizado (seis dígitos centrais). \\
captura\_id & texto(32) & não & FK capturas & Captura que originou a entidade (chave estrangeira para capturas); as demais ficam em entidade\_fontes. \\
\hline
\end{longtable}
\end{small}

\section*{extracoes\_ia}
Execução da extração assistida por IA sobre uma captura, com os dados para reconstituí-la e os itens recusados (RF13, RNF05).

\begin{small}
\begin{longtable}{p{3.4cm}p{1.9cm}p{0.8cm}p{1.6cm}p{6.6cm}}
\hline
\textbf{Coluna} & \textbf{Tipo} & \textbf{Nulo} & \textbf{Chave} & \textbf{Descrição} \\ \hline\endhead
id & texto(32) & não & PK & Identificador da execução, UUID4 em hexadecimal. \\
caso\_id & texto(32) & não & FK casos & Caso da captura analisada (chave estrangeira para casos). \\
captura\_id & texto(32) & não & FK capturas & Captura analisada (chave estrangeira para capturas). \\
captura\_sha256 & texto(64) & não &  & SHA-256 da captura analisada, copiado no início da execução. \\
provedor & texto(30) & não &  & Provedor de IA usado. Valores no código: ollama, falso (provedor de teste). \\
modelo & texto(200) & não &  & Nome do modelo informado pelo provedor; vazio se a execução falhou antes de identificá-lo. \\
modelo\_versao & texto(200) & não &  & Versão do modelo (resumo do modelo) informada pelo provedor. \\
servidor\_versao & texto(50) & não &  & Versão do servidor do provedor, informada por ele. \\
parametros & JSON & não &  & JSON com os parâmetros do provedor, tamanho do bloco, limite de caracteres, se o texto foi truncado e a codificação. \\
prompt\_sha256 & texto(64) & não &  & SHA-256 das instruções enviadas ao modelo somadas ao esquema JSON da resposta; constante no código. \\
situacao & texto(15) & não &  & Situação da execução. Valores: em\_andamento, sucesso, falha. \\
motivo & texto & não &  & Motivo da falha, até 400 caracteres; vazio em caso de sucesso. \\
blocos & inteiro & não &  & Número de blocos em que o texto da captura foi dividido. \\
itens\_recebidos & inteiro & não &  & Quantidade de itens devolvidos pelo modelo. \\
itens\_aceitos & inteiro & não &  & Quantidade de itens que passaram nas verificações e viraram sugestões em sugestoes\_ia. \\
itens\_recusados & inteiro & não &  & Quantidade de itens recusados nas verificações. \\
recusas & JSON & não &  & Lista JSON, até 200 itens, com o motivo da recusa (por exemplo trecho\_nao\_literal, valor\_fora\_do\_trecho) e resumo do item. \\
usuario & texto(100) & não &  & Login de quem pediu a extração. \\
iniciado\_em & texto(40) & não &  & Data e hora de início, ISO 8601 em UTC. \\
concluido\_em & texto(40) & sim &  & Data e hora de término, ISO 8601 em UTC; nulo enquanto em andamento. \\
\hline
\end{longtable}
\end{small}

\section*{arestas}
Vínculo dirigido entre duas entidades do caso, derivado das capturas; forma o grafo de relacionamentos (RF02).

\begin{small}
\begin{longtable}{p{3.4cm}p{1.9cm}p{0.8cm}p{1.6cm}p{6.6cm}}
\hline
\textbf{Coluna} & \textbf{Tipo} & \textbf{Nulo} & \textbf{Chave} & \textbf{Descrição} \\ \hline\endhead
id & texto(32) & não & PK & Identificador da aresta, UUID4 em hexadecimal. \\
caso\_id & texto(32) & não & FK casos & Caso a que a aresta pertence (chave estrangeira para casos). \\
origem\_id & texto(32) & não & FK entidades & Entidade de origem do vínculo (chave estrangeira para entidades). \\
destino\_id & texto(32) & não & FK entidades & Entidade de destino do vínculo (chave estrangeira para entidades). \\
tipo & texto(30) & não &  & Tipo do vínculo. Valores no código: socio\_de, contratado\_por, fornecedor\_de, contratante\_de, mesmo\_endereco, sancionada\_em, servidor\_de, doou\_para, edital\_de, citado\_em. \\
chave & texto(200) & não &  & Discriminador do vínculo, como número de controle do contrato ou endereço normalizado; pode ficar vazio. \\
atributos & JSON & não &  & JSON com dados do vínculo; em citado\_em, o trecho, a posição e a sugestão de IA de origem. \\
captura\_id & texto(32) & não & FK capturas & Captura que fundamenta o vínculo (chave estrangeira para capturas). \\
\hline
\end{longtable}
\end{small}

\section*{entidade\_fontes}
Relação entre cada entidade e todas as capturas que contribuíram para ela; registra a proveniência dos dados da entidade (cadeia de custódia, RF07 no README).

\begin{small}
\begin{longtable}{p{3.4cm}p{1.9cm}p{0.8cm}p{1.6cm}p{6.6cm}}
\hline
\textbf{Coluna} & \textbf{Tipo} & \textbf{Nulo} & \textbf{Chave} & \textbf{Descrição} \\ \hline\endhead
id & inteiro & não & PK & Identificador numérico sequencial, gerado pelo banco. \\
entidade\_id & texto(32) & não & FK entidades & Entidade (chave estrangeira para entidades). \\
captura\_id & texto(32) & não & FK capturas & Captura que contribuiu para a entidade (chave estrangeira para capturas). \\
papel & texto(50) & não &  & Papel da captura em relação à entidade, definido pelo resolvedor, por exemplo orgaoEntidade, contrato, edital, orgao\_servidores, extracao\_ia. \\
\hline
\end{longtable}
\end{small}

\section*{monitoramentos}
Entidade marcada para acompanhamento contínuo; o agendador refaz as coletas de origem na periodicidade escolhida, no máximo um por entidade (RF06).

\begin{small}
\begin{longtable}{p{3.4cm}p{1.9cm}p{0.8cm}p{1.6cm}p{6.6cm}}
\hline
\textbf{Coluna} & \textbf{Tipo} & \textbf{Nulo} & \textbf{Chave} & \textbf{Descrição} \\ \hline\endhead
id & texto(32) & não & PK & Identificador do monitoramento, UUID4 em hexadecimal. \\
caso\_id & texto(32) & não & FK casos & Caso do monitoramento (chave estrangeira para casos). \\
entidade\_id & texto(32) & não & FK entidades & Entidade monitorada (chave estrangeira para entidades). \\
periodicidade\_horas & inteiro & não &  & Intervalo entre execuções, em horas; o mínimo vem da configuração (padrão 1 hora). \\
ativo & lógico & não &  & Indica se o agendador deve executar o monitoramento. \\
motivo\_inativo & texto & não &  & Por que foi desativado: pelo usuário ('Desativado por <login>.'), prazo de retenção vencido ou entidade inexistente; vazio quando ativo. \\
justificativa & texto & não &  & Justificativa de quem marcou a entidade; acompanha cada execução automática na auditoria. \\
criado\_por & texto(100) & não &  & Login de quem criou o monitoramento. \\
criado\_em & texto(40) & não &  & Data e hora de criação, ISO 8601 em UTC. \\
proxima\_execucao\_em & texto(40) & não &  & Data e hora da próxima execução prevista, ISO 8601 em UTC. \\
ultima\_execucao\_em & texto(40) & sim &  & Data e hora do fim da última execução, ISO 8601 em UTC; nulo se nunca executou. \\
\hline
\end{longtable}
\end{small}

\section*{sugestoes}
Sugestão de que duas entidades do caso podem ser a mesma pessoa (homônimos); o sistema nunca as funde automaticamente (RF02).

\begin{small}
\begin{longtable}{p{3.4cm}p{1.9cm}p{0.8cm}p{1.6cm}p{6.6cm}}
\hline
\textbf{Coluna} & \textbf{Tipo} & \textbf{Nulo} & \textbf{Chave} & \textbf{Descrição} \\ \hline\endhead
id & texto(32) & não & PK & Identificador da sugestão, UUID4 em hexadecimal. \\
caso\_id & texto(32) & não & FK casos & Caso das entidades (chave estrangeira para casos). \\
entidade\_id & texto(32) & não & FK entidades & Entidade de referência (chave estrangeira para entidades). \\
candidata\_id & texto(32) & não & FK entidades & Entidade candidata a ser a mesma (chave estrangeira para entidades); cada par é gravado nos dois sentidos. \\
score & real & não &  & Grau de semelhança. Valores no código: 0.9 (mesmo nome e mesmo CPF descaracterizado), 0.7 (CPF ausente em um), 0.5 (CPFs diferentes). \\
motivo & texto & não &  & Texto gerado pelo resolvedor explicando a semelhança e pedindo confirmação antes de unificar. \\
criado\_em & texto(40) & não &  & Data e hora da sugestão, ISO 8601 em UTC. \\
\hline
\end{longtable}
\end{small}

\section*{sugestoes\_ia}
Item extraído pela IA à espera da decisão do investigador; só vira entidade se aceito (RF13, RNF05).

\begin{small}
\begin{longtable}{p{3.4cm}p{1.9cm}p{0.8cm}p{1.6cm}p{6.6cm}}
\hline
\textbf{Coluna} & \textbf{Tipo} & \textbf{Nulo} & \textbf{Chave} & \textbf{Descrição} \\ \hline\endhead
id & texto(32) & não & PK & Identificador da sugestão, UUID4 em hexadecimal. \\
caso\_id & texto(32) & não & FK casos & Caso da captura (chave estrangeira para casos). \\
extracao\_id & texto(32) & não & FK extracoes\_ia & Execução que gerou a sugestão (chave estrangeira para extracoes\_ia). \\
captura\_id & texto(32) & não & FK capturas & Captura de onde o item foi extraído (chave estrangeira para capturas). \\
tipo & texto(20) & não &  & Tipo do item. Valores: empresa, orgao, pessoa, valor, data, objeto. \\
valor & texto(500) & não &  & Valor extraído, como escrito no documento, até 500 caracteres. \\
campos & JSON & não &  & JSON com campos opcionais cnpj, cpf e papel; o CPF é gravado descaracterizado. \\
trecho & texto & não &  & Citação literal do texto da captura que contém o valor; verificada por comparação exata. \\
posicao\_inicio & inteiro & não &  & Posição inicial do trecho no texto da captura, em caracteres. \\
posicao\_fim & inteiro & não &  & Posição final do trecho no texto da captura, em caracteres. \\
confianca & real & sim &  & Confiança informada pelo modelo, entre 0 e 1; nula se o modelo não informou. \\
situacao & texto(10) & não &  & Situação da sugestão. Valores: pendente, aceita, recusada. \\
decidido\_por & texto(100) & sim &  & Login de quem decidiu; nulo enquanto pendente. \\
decidido\_em & texto(40) & sim &  & Data e hora da decisão, ISO 8601 em UTC; nulo enquanto pendente. \\
justificativa\_decisao & texto & sim &  & Justificativa da decisão; nula enquanto pendente. \\
entidade\_id & texto(32) & sim & FK entidades & Entidade criada ou reaproveitada ao aceitar (chave estrangeira para entidades); nula se recusada ou se o tipo não gera entidade. \\
\hline
\end{longtable}
\end{small}

\section*{monitoramento\_execucoes}
Cada execução de um monitoramento, automática ou manual, com as recoletas feitas e a quantidade de avisos gerados (RF06).

\begin{small}
\begin{longtable}{p{3.4cm}p{1.9cm}p{0.8cm}p{1.6cm}p{6.6cm}}
\hline
\textbf{Coluna} & \textbf{Tipo} & \textbf{Nulo} & \textbf{Chave} & \textbf{Descrição} \\ \hline\endhead
id & texto(32) & não & PK & Identificador da execução, UUID4 em hexadecimal. \\
monitoramento\_id & texto(32) & não & FK monitoramentos & Monitoramento executado (chave estrangeira para monitoramentos). \\
caso\_id & texto(32) & não & FK casos & Caso do monitoramento (chave estrangeira para casos). \\
origem & texto(12) & não &  & Quem disparou a execução. Valores: agendador, manual. \\
usuario & texto(100) & não &  & Login de quem executou; 'agendador' nas execuções automáticas. \\
iniciado\_em & texto(40) & não &  & Data e hora de início, ISO 8601 em UTC. \\
concluido\_em & texto(40) & sim &  & Data e hora de término, ISO 8601 em UTC; nulo enquanto em andamento. \\
situacao & texto(20) & não &  & Resultado da execução. Valores: em\_andamento, sucesso, parcial, falha, sem\_coletas, recusada\_retencao. \\
coletas & JSON & não &  & Lista JSON, uma linha por recoleta: coleta\_id nova, recoleta\_de, fonte, situação, motivo e mesmo\_conteudo (IDs apontam para coletas, sem chave estrangeira). \\
avisos\_gerados & inteiro & não &  & Quantidade de avisos criados nesta execução. \\
motivo & texto & não &  & Explicação da situação quando não é sucesso; vazio em caso de sucesso. \\
\hline
\end{longtable}
\end{small}

\section*{avisos}
Novidade detectada pelo monitoramento sobre uma entidade, como mudança de sócio, contrato, sanção ou termo, com controle de leitura (RF06).

\begin{small}
\begin{longtable}{p{3.4cm}p{1.9cm}p{0.8cm}p{1.6cm}p{6.6cm}}
\hline
\textbf{Coluna} & \textbf{Tipo} & \textbf{Nulo} & \textbf{Chave} & \textbf{Descrição} \\ \hline\endhead
id & texto(32) & não & PK & Identificador do aviso, UUID4 em hexadecimal. \\
caso\_id & texto(32) & não & FK casos & Caso do aviso (chave estrangeira para casos). \\
monitoramento\_id & texto(32) & não & FK monitoramentos & Monitoramento que gerou o aviso (chave estrangeira para monitoramentos). \\
execucao\_id & texto(32) & não & FK monitoramento\_execucoes & Execução em que a novidade foi detectada (chave estrangeira para monitoramento\_execucoes). \\
entidade\_id & texto(32) & não & FK entidades & Entidade a que a novidade se refere (chave estrangeira para entidades). \\
tipo & texto(30) & não &  & Tipo da novidade. Valores: quadro\_societario, contrato\_novo, sancao\_nova, termo\_novo. \\
titulo & texto(300) & não &  & Título do aviso, gerado pelo monitoramento. \\
descricao & texto & não &  & Descrição da novidade, gerada pelo monitoramento. \\
detalhes & JSON & não &  & JSON com os dados da novidade, como o registro e o tipo de mudança. \\
captura\_ids & JSON & não &  & Lista JSON de IDs das capturas em que a novidade aparece (aponta para capturas, sem chave estrangeira). \\
criado\_em & texto(40) & não &  & Data e hora de criação do aviso, ISO 8601 em UTC. \\
lido\_em & texto(40) & sim &  & Data e hora em que o aviso foi marcado como lido, ISO 8601 em UTC; nulo se não lido. \\
lido\_por & texto(100) & sim &  & Login de quem marcou o aviso como lido; nulo se não lido. \\
\hline
\end{longtable}
\end{small}

